\ficha{Eichengreen (1996), Globalizing Capital}{eichengreen1996}

\begin{identificacao}
EICHENGREEN, Barry. \emph{Globalizing Capital}: a history of the international
monetary system. Princeton: Princeton University Press, 1996.

Livro de síntese, originado de quatro conferências Gaston Eyskens na
Universidade Católica de Leuven (p. vii). Leitura integral do prefácio
(pp. vii-viii), do capítulo 1, \emph{Introduction} (pp. 3-6), e do capítulo 2,
\emph{The Gold Standard} (pp. 7-44); leitura por alto do capítulo 6,
\emph{Conclusion} (pp. 188-192); glossário (pp. 193-199) consultado. Os
capítulos 3 a 5, sobre o entreguerras, Bretton Woods e a flutuação, ficaram
fora do recorte. Idioma original, inglês. Edição lida, a primeira, de 1996.
\end{identificacao}

\bloco{Problema e tese}
O livro pergunta por que taxas de câmbio fixas conviveram com alta mobilidade
de capital antes de 1914 e deixaram de conviver depois de 1945. A resposta é
política, não técnica: o que sustenta uma paridade é a proteção do governo
contra a pressão para trocar estabilidade cambial por outros objetivos. Antes
de 1914 essa proteção vinha do insulamento em relação à política doméstica, com
sufrágio restrito, sindicatos fracos e partidos trabalhistas incipientes; depois
da guerra veio dos controles de capital; quando os dois se esgotaram, restou
flutuar ou unificar a moeda (p. 4, p. 191).

\bloco{Argumento}
O livro inteiro organiza a história monetária internacional de 1850 aos anos
1990 em quatro eras do mercado de capitais, cuja mobilidade desenha um U ao
longo do tempo (p. 3). Esse desenho derruba a explicação corrente para o fim de
Bretton Woods, segundo a qual a flutuação seria consequência inevitável da
mobilidade de capital: a mobilidade também era alta antes de 1914 e não impediu
a paridade fixa (p. 4). A variável que muda é a política interna, e os controles
de capital do pós-guerra foram um substituto funcional do sufrágio restrito do
século XIX (p. 5). Externalidades de rede e dependência de trajetória explicam
por que cada país adota o padrão de seus parceiros e por que reformar o sistema
é empreendimento coletivo, fracassado nas conferências de 1870, 1920 e 1970
(p. 6).

O recorte lido sustenta esse argumento em oito passos.

\begin{enumerate}[leftmargin=*,nosep]
\item O padrão-ouro clássico não foi negociado. Emergiu de escolhas nacionais
constrangidas pela decisão dos vizinhos, a partir da adoção britânica acidental
de 1717, quando Newton fixou preço baixo demais para a prata (p. 7).
\item O bimetalismo durou até 1870 não pelo lobby da prata nem por limite
técnico de cunhagem, mas porque externalidades de rede o mantinham no lugar
(p. 15). Flandreau não encontra nas comissões monetárias europeias a clivagem
uniforme entre agricultores e industriais (p. 15). A reação em cadeia dos anos
1870 veio do incentivo de cada país a adotar o padrão de seus vizinhos
comerciais e financeiros, não da liquidação alemã de prata (p. 18).
\item O sistema nunca foi homogêneo. Apenas Inglaterra, Alemanha, França e
Estados Unidos mantinham padrão-ouro puro (p. 20); os demais operavam com papel,
prata e moeda fiduciária, sob sistemas fiduciários ou proporcionais, com
\emph{gold devices} e reservas acima do mínimo legal (pp. 21-24). Índia,
Filipinas e boa parte da América Latina guardavam reservas inteiramente em
divisas (p. 23).
\item A conduta dos bancos centrais não obedecia a nenhuma regra escrita. A
expressão ``regras do jogo'' é de Keynes, em 1925 (p. 28), e Bloomfield mostrou
que as violações eram tão frequentes antes de 1913 quanto Nurkse as encontrara
no entreguerras (pp. 28-29). Lucratividade, custo da dívida pública e efeito
sobre a atividade influíam na taxa de desconto (p. 29).
\item A explicação do funcionamento é a credibilidade. O padrão-ouro é tratado
como instituição historicamente específica.
\chave{It was a socially constructed institution whose viability hinged on the
context in which it operated.}{p. 30}
A prioridade dada à conversibilidade não era contestada porque o voto era
limitado a homens de propriedade, os partidos trabalhistas estavam em formação e
salários e preços eram flexíveis (p. 31). Investidores sabiam disso, e o capital
fluía em direção estabilizadora quando o câmbio se aproximava do ponto de
exportação do ouro (p. 31).
\chave{Central banks possessed the capacity to violate the rules of the game in
the short run because there was no question about obeying them in the long
run.}{p. 32}
\item A credibilidade era complementada por cooperação. Na crise Baring, em
1890, o Banco da Inglaterra tomou emprestadas 3 milhões de libras em ouro do
Banco da França e obteve promessa de 1,5 milhão da Rússia (p. 34). Seguiram-se
episódios em 1893, 1898, 1906, 1907, 1909 e 1910, e o autor conclui que a
sobrevivência do sistema dependia da colaboração entre bancos centrais e
governos, não de automatismo (p. 35). Cláusulas de escape permitiam suspender a
conversibilidade em circunstância verificável e alheia à vontade das autoridades
sem destruir a credibilidade do compromisso em tempos normais (p. 37).
\item Na periferia nada disso se reproduzia. A cooperação raramente chegava lá,
porque problemas periféricos não ameaçavam a estabilidade sistêmica (p. 38).
Faltavam bancos centrais, criados na América Latina só nos anos 1920 (p. 39). A
pauta de exportação concentrada expunha esses países a choques de termos de
troca, e choques de conta corrente e de conta de capital se reforçavam, em vez
de se compensarem como no caso britânico (p. 39). Argentina, Brasil, Chile,
Itália e Portugal suspenderam repetidamente a conversibilidade (p. 41), e a
depreciação era bem recebida por proprietários endividados e por exportadores,
frequentemente as mesmas pessoas (p. 41).
\chave{Over much of the world, the absence of the special political and social
factors that lent the gold standard its credibility at the system's European
core rendered its operation problematic.}{p. 42}
\item A estabilidade do sistema foi exceção histórica, não norma (p. 42), e já
se erodia antes de 1914: a posição singular da Grã-Bretanha declinava, a
adequação do estoque de ouro voltava a preocupar, os Estados Unidos eram fonte
recorrente de choque, o sufrágio se estendia e as tensões posteriores à partilha
da África corroíam a solidariedade entre bancos centrais (pp. 42-43).
\end{enumerate}

\bloco{Conceitos e definições operacionais}
Padrão-ouro é definido pela prática, não pelo estatuto: um país está no padrão
quando o governo se compromete a converter sua moeda em ouro a preço fixo e
mantém reserva para isso, qualquer que seja o meio circulante interno (p. 22).
Conversibilidade, no glossário, é a troca livre da moeda por ouro a preço fixo
(p. 194). Credibilidade não é definida formalmente e opera como ausência de
dúvida, entre quem tem poder de influir na política, de que as autoridades farão
o necessário para defender a paridade (p. 30). Regra é o oposto de código
escrito: não existe manual de conduta, e o que existe é consenso político
(p. 30). Cláusula de escape é a abrogação temporária de uma regra de política
econômica (p. 194), válida quando a contingência é verificável e não produzida
pelas próprias autoridades (p. 37). Padrão ouro-câmbio é o sistema em que as
reservas podem tomar a forma de moedas conversíveis, e não só de ouro (p. 196).
\emph{Currency board} aparece apenas no glossário, como substituto do banco
central que amarra a política monetária à de outro país por lei ou constituição
(p. 194).

\bloco{Evidência e método}
Síntese interpretativa da literatura secundária, sem dado novo. O recorte cobre
de 1717 a 1914 e apoia-se em quatro peças: preço relativo do ouro e da prata,
1830-1902, de Warren e Pearson (Figura 2.1); preços por atacado britânicos,
1873-1913, de Mitchell (Figura 2.2); a estrutura do padrão-ouro pós-1880 por
forma de circulação e de reserva (Figura 2.3); e a composição das reservas em
divisas, 1899-1913, de Lindert (Tabela 2.1). A prova central sobre as regras do
jogo é a replicação feita por Bloomfield (1959) do exercício de Nurkse (1944),
que contava a correlação entre ativos internos e externos dos bancos centrais
(pp. 28-29). Há dois contrafactuais discutidos: o de Friedman, segundo o qual o
bimetalismo internacional teria dado nível de preços mais estável (p. 19), e o
cálculo de Oppers, segundo o qual a desmonetização alemã não teria rompido a
razão de 15,5 para 1 nos países da União Latina (p. 18). Não há econometria
própria no recorte.

\bloco{Agentes e vertentes}
Teoria e modelo: Hume e o modelo de fluxo de preços e espécie (p. 25); o Comitê
Cunliffe (p. 26); Keynes, autor da expressão sobre as regras do jogo e da imagem
do Banco da Inglaterra como regente da orquestra internacional (pp. 28 e 33);
Bagehot e a regra de descontar livremente na drenagem interna (p. 37). Leitura
empírica do sistema: Nurkse, Bloomfield, Whale, Machlup, Ohlin e Lindert.
Interpretação política: Polanyi, cuja tese sobre a reação social à expansão do
mercado o livro se propõe a reexaminar (p. 5). Historiografia do bimetalismo:
Redish, com o argumento técnico da prensa a vapor; Flandreau e Frieden, contra a
clivagem simples entre agrários e industriais; Gallarotti e Oppers. Sobre
cláusula de escape, Bordo e Kydland, além do próprio autor (p. 37). Agentes
históricos: Newton, Napoleão, Bryan e McKinley na disputa americana da prata
(p. 41), o Banco da Inglaterra, o Banco da França, o Reichsbank, o Banco Estatal
Russo e a casa Baring Brothers.

\bloco{Críticas e limites}
A tese política é afirmada, não medida: no recorte não há teste da relação entre
extensão do sufrágio e defesa da paridade, e o mecanismo se apoia em
plausibilidade e em citação de contemporâneos. O contraste entre centro e
periferia é construído sobretudo com o caso dos Estados Unidos, ao qual o autor
dedica três páginas (pp. 39-41), enquanto a América Latina ocupa pouco mais de
uma. O Brasil aparece duas vezes no recorte e sem análise: como fonte do ouro
que desmonetizou a prata inglesa no século XVIII (p. 12) e num rol de países que
suspenderam repetidamente a conversibilidade (p. 41). O texto também não
hierarquiza os fatores da instabilidade periférica, que misturam estrutura de
exportação, ausência de banco central e coalizões pró-depreciação. E não separa
insulamento político interno de credibilidade perante o credor externo,
distinção que a literatura posterior sobre sinalização, como Bordo e Rockoff e
Obstfeld e Taylor, tomou como o problema central da periferia. Por fim, o
recorte não trata de caixas de conversão, e a única discussão de \emph{currency
board} está no capítulo 5, sobre a Argentina de 1991, fora desta leitura.

\begin{dialogo}
\item Procurar nos jornais o argumento de que a Caixa de Conversão vale por
submeter a emissão a uma regra automática, e a réplica de que o que decide é o
nível da taxa. Eichengreen sustenta que a regra nunca foi mecânica e que o que
sustentava a paridade era consenso político (p. 30). A ponte testa se o
vocabulário brasileiro de garantia é o do compromisso incontestado ou o da
defesa de interesse exportador.
\item Procurar menções a Londres, aos Rothschild, ao saldo em ouro depositado
fora e à agência de Londres da Caixa. O texto afirma que os governos que tomavam
emprestado em Londres, Paris ou Berlim eram obrigados pelos credores a manter
parte dos recursos depositada na praça, e que mesmo sem exigência mantinham o
depósito como sinal de crédito (p. 23). Hipótese a testar: a agência de Londres
é apresentada na imprensa como sinal de credibilidade, e não apenas como
conveniência operacional.
\item Procurar, no debate de 1914, a justificativa da suspensão do troco. Se a
imprensa a apresenta como resposta a circunstância externa verificável e alheia
ao governo, o caso brasileiro se encaixa na cláusula de escape descrita na
p. 37; se a apresenta como fracasso do plano de 1906, não se encaixa.
\item Procurar quem defende a taxa baixa e a que interesse se associa. O modelo
prevê coalizão de proprietários endividados com exportadores, muitas vezes as
mesmas pessoas (p. 41). Testar contra o eixo valorização e emissão nos jornais
paulistas, verificando se lavoura e comércio exportador aparecem alinhados como
o modelo supõe.
\item Procurar se a imprensa descreve a Caixa como entrada num clube
internacional. Eichengreen lista a difusão do padrão entre 1890 e 1900 sem
incluir o Brasil (p. 18) e o situa entre os países que suspendiam com frequência
(p. 41). A ponte testa se o discurso de 1906 reivindica pertencimento ao sistema
ou apenas estabilidade interna do câmbio.
\end{dialogo}

\usonocapitulo{Parte I. Fornece o enquadramento geral da seção: o padrão-ouro
como instituição historicamente específica, cuja viabilidade dependia de
credibilidade e de cooperação entre bancos centrais, e não de automatismo; e a
assimetria entre centro e periferia como consequência da ausência, fora da
Europa, do arranjo político que dava crédito ao compromisso. Serve também de
contraponto a Bordo e Kydland e a Bordo e Rockoff, porque desloca a explicação
da sinalização ao mercado de capitais para o insulamento político e a
solidariedade entre bancos centrais. Não serve à Parte IV: o recorte não trata
de caixas de conversão, e a passagem sobre a América Latina é curta demais para
sustentar afirmação sobre o Brasil, que precisa vir de Fritsch e Franco, Topik e
Franco.}
